\chapter{Results}
\label{ch:results}

\begin{guidance}[What this chapter does]
Presents the results objectively, without yet judging them (that is
\cref{ch:evaluation}). Use figures and tables, label axes with quantities and
units, and refer to every figure in the text. State uncertainties. Typical
length: 800--1,200 words.

Plots are best made directly from your data with \texttt{pgfplots}, as in
\cref{fig:cooling}: the figure then updates when the CSV file in
\texttt{data/} changes, and the fonts match the report. Exported images
(PDF or PNG from MATLAB, Python, Excel) are also fine; save them in
\texttt{figures/}.
\end{guidance}

\section{Temperature accuracy (test T1)}
\label{sec:res-accuracy}

\Cref{fig:cooling} compares the logged temperature with the reference
during a cooling test. The \gls{mae} over the test was \SI{0.21}{\celsius},
within the limit of \cref{req:accuracy}.

\begin{figure}[htbp]
  \centering
  \begin{tikzpicture}
    \begin{axis}[
        width=0.85\linewidth, height=6cm,
        xlabel={Time $t$ in \si{\minute}},
        ylabel={Temperature $T$ in \si{\celsius}},
        grid=major, grid style={black!10},
        legend pos=north east, legend cell align=left,
        every axis plot/.append style={thick}]
      \addplot[SeaBlue, mark=*, mark size=1.2pt]
        table[col sep=comma, x=time_min, y=logger_C] {data/cooling-test.csv};
      \addlegendentry{Data logger}
      \addplot[SolentRed, dashed]
        table[col sep=comma, x=time_min, y=reference_C] {data/cooling-test.csv};
      \addlegendentry{Calibrated reference}
    \end{axis}
  \end{tikzpicture}
  \caption{Logged and reference temperature during the cooling test T1 (example data)}
  \label{fig:cooling}
\end{figure}

\section{Battery life (test T3)}
\label{sec:res-battery}
